% ============================================================================================
% BAB VI EVALUASI — KERANGKA (bullet), bukan prosa jadi.
%
% Menampung DECISIONS.md "Bab IV — Hasil dan Pembahasan". Lihat catatan pemetaan
% di kepala berkas "Bab III - Analisis.tex": "Bab IV" pada DECISIONS.md berarti
% HASIL, dan berlabuh di sini.
%
% ATURAN YANG DIPATUHI BERKAS INI: TIDAK ADA ANGKA HASIL YANG DIKARANG.
% Keluaran run final BELUM DIBACA. Karena itu:
%   - tabel dan gambar DISEBUTKAN NAMANYA beserta kolom/sumbunya, tetapi belum
%     ditulis sebagai float. Menuliskannya sebagai \begin{table} berisi angka
%     penampung akan memasukkan entri palsu ke Daftar Tabel dan Daftar Gambar,
%     dan angka penampung punya kebiasaan bertahan sampai cetak;
%   - angka yang MUNCUL di bawah ini seluruhnya berasal dari DECISIONS.md sebagai
%     hasil TERUKUR yang sudah tercatat (D-14, D-27, D-20, D-04, D-03), bukan
%     dari run final. Provenance-nya ditandai di tempat.
%
% PENYIMPANGAN DARI TEMPLATE (D-45). Subbab template "Evaluasi Kesalahan
% Penulisan yang Sering Terjadi" beserta enam anak-subbabnya adalah instruksi
% gaya penulisan, bukan isi tesis, sehingga dihapus. Aturannya tetap berlaku dan
% ada pada Template Baskara serta pada skill itb-sti-thesis.
%
% Tidak ada \autocite / \cite di berkas ini.
% ============================================================================================
\cleardoublepage
\chapter{EVALUASI}
\label{chap:evaluasi}

\section{Metode Evaluasi}

\begin{itemize}
    \item \textbf{Satuan pelaporan.} Seluruh metrik dilaporkan dalam satuan GFD
          (sambaran per km\textsuperscript{2} per tahun) setelah inversi dan
          agregasi, mengikuti urutan yang dikunci D-03: \textit{un-standardise}
          $\rightarrow$ \texttt{expm1} $\rightarrow$ jumlahkan.
    \item \textbf{Metrik regresi:} RMSE, MAE, R\textsuperscript{2}, NRMSE, dan
          \textbf{bias} $=$ rerata(prediksi) $-$ rerata(observasi). NRMSE
          dinormalkan terhadap \textbf{julat} teramati --- konvensi yang dipakai
          saat NF-01 ditulis --- sehingga dilaporkan di samping RMSE, bukan
          menggantikannya. Bias lebih penting di sini daripada biasanya: D-04
          menjumlahkan prediksi per jam menjadi jendela, dan bias per jam yang
          sistematis \textbf{tidak saling meniadakan di bawah penjumlahan,
          melainkan berlipat}.
    \item \textbf{Metrik klasifikasi} (tahap okurensi): \textit{log loss},
          AUC, \textit{average precision}, Brier, dan F1 pada satu ambang.
          AUC dan \textit{average precision} bebas ambang, yang penting pada
          laju dasar 5,7\%; \textbf{\textit{average precision} adalah headline
          yang lebih jujur}, karena AUC menyanjung masalah dengan positif jarang.
          Ambang bawaan adalah \textbf{laju dasar teramati}, bukan 0,5.
    \item \textbf{Lantai (\textit{baseline}) menyertai setiap tabel.} Setiap
          tabel metrik memuat skor \textit{baseline} trivial dalam satuan yang
          sama, karena R\textsuperscript{2} tanpa lantai bukan sekadar tidak
          lengkap melainkan tidak dapat ditafsirkan (D-07, D-12, D-20).
          \texttt{skill\_score} $= 1 - \text{RMSE}/\text{RMSE}_{\text{baseline}}$.
    \item \textbf{Faktor \textit{smearing}} (D-03): koreksi multiplikatif satu
          parameter atas bias retransformasi, versi satu parameter dari
          \textit{estimator} Duan. \textbf{Dipasang pada validasi, tidak pernah
          pada uji} --- faktor yang dipasang pada baris uji adalah kebocoran
          target yang menyamar sebagai kalibrasi.
          % CITE: Duan, smearing estimator
    \item \textbf{Kalibrasi per jendela}: rerata prediksi terhadap rerata
          observasi, dengan \texttt{ratio} $<1$ berarti prediksi terlalu rendah.
          Dilaporkan \textbf{di samping setiap hasil berjendela}.
    \item \textbf{Replikasi.} Tabel utama: 3 fold (D-06) $\times$ 3 seed (D-29),
          dilaporkan rerata $\pm$ simpangan baku. \textit{Sweep}: satu fold, satu
          seed (D-37). \textbf{Asimetri ini dinyatakan di sini, bukan ditemukan
          pembaca.}
\end{itemize}

\section{Hasil Evaluasi}

\subsection{Batasan yang Membingkai Setiap Hasil}

% Subbab ini mendahului hasil, bukan mengikutinya. Setiap butir di bawah adalah
% hal yang harus dinyatakan SEBELUM angka dibaca.

\begin{itemize}
    \item \textbf{Seluruh hasil subtropis bersifat provisional.}
          % TODO(open): pertanyaan terbuka 10 — komposisi CG/IC MERLIN. Selama
          %   belum diselesaikan lewat dokumentasi KSC untuk produk
          %   MerlinCloudToGround, setiap angka subtropis di bab ini harus
          %   membawa pernyataan itu. Rinciannya di Bab III.
    \item \textbf{Gradien \textit{detection efficiency} subtropis} (D-27):
          Spearman $-0,955$ terhadap jarak dari Cape, melintasi empat orde
          besaran [terukur, tercatat --- bukan dari run final]. Dilaporkan
          sebagai konfound antara respons instrumen dan klimatologi yang tidak
          dapat dipisahkan dengan MERLIN saja.
    \item \textbf{Distribusi \texttt{coverage} atas baris yang benar-benar
          dilatih} (D-10, D-18), dilaporkan di samping pangsa nol.
          \texttt{coverage} \textbf{tidak pernah} dipakai sebagai filter maupun
          bobot sampel, karena ia berkorelasi dengan target melalui aktivitas
          petir itu sendiri --- bulan yang benar-benar tenang dan detektor yang
          mati terlihat identik.
          % TODO: nyatakan PEMBOBOTAN angka cakupan — berbobot BARIS atau
          %   berbobot BULAN. Angka berbobot baris: tropis rerata 0,8549 min
          %   0,1290; subtropis rerata 0,5921 min 0,0323. Versi lama menulis
          %   tropis 86%, yang kemungkinan berbobot bulan.
    \item \textbf{Pangsa nol di samping setiap metrik} (D-20): 94,30\% tropis,
          97,00\% subtropis pada resolusi jam.
    \item \textbf{Asimetri cakupan antar domain} (D-26): 84 dari 84 bulan tropis
          lawan 81 dari 84 subtropis; tiga bulan absen, bukan diisi nol.
    \item \textbf{Siaran AOD} (D-24): batasan cara parameter harian/bulanan
          disebarkan ke baris per jam --- nyatakan bagaimana, karena ia
          memengaruhi tafsir kepentingan fitur.
    \item \textbf{Sel semua-nol di bawah pemisahan sel} (D-11): enam sel tropis,
          20\% dari grid, tanpa satu pun sambaran sepanjang rekaman.
    \item \textbf{Tujuh hiperparameter tanpa alasan tercatat} (pertanyaan
          terbuka 22) --- keterbatasan yang dinyatakan, bukan lubang yang diisi
          dengan pembenaran yang ditulis sekarang.
\end{itemize}

\subsection{Hasil Tangga Model}

\begin{itemize}
    \item \textbf{Tabel: tangga enam \textit{rung} per domain per tahap}
          (D-34). Baris: \texttt{baseline\_trivial}, \texttt{ridge},
          \texttt{nn\_matched}, \texttt{qnn}, \texttt{nn\_large},
          \texttt{nn\_full}. Kolom: jumlah parameter, metrik tahap okurensi
          (\textit{log loss}, AP, AUC, Brier), metrik tahap cacah (RMSE, MAE,
          R\textsuperscript{2}, bias), dan \texttt{skill\_score}. Nilai sebagai
          rerata $\pm$ simpangan baku atas 3 fold $\times$ 3 seed.
          % TODO: isi setelah keluaran run final dibaca.
    \item \textbf{Paritas parameter 47 lawan 52} dinyatakan di badan tabel, dan
          \texttt{nn\_full} ditandai sebagai \textbf{sengaja mematahkan paritas}.
    \item \textbf{\texttt{nn\_full} di bawah anggaran langkah yang disamakan}
          (D-30): anggaran sama, data lebih banyak. Inilah pembacaan yang benar
          atas \textit{rung} tersebut, dan ia berbeda dari versi pra-perbaikan
          yang menghasilkan sesuatu yang \textbf{tampak} seperti temuan.
          % TODO: isi setelah run final dibaca.
\end{itemize}

\subsection{Kurva Jendela Pelaporan}

\begin{itemize}
    \item \textbf{Gambar: kurva keterampilan terhadap resolusi pelaporan}
          (D-04), satu titik per jendela 1/3/6/24 jam, per domain, per model.
    \item \textbf{Tabel: metrik per jendela}, membawa pangsa nol teramati per
          jendela agar pembaca melihat target berubah bentuk saat jendela
          melebar --- 94,30/97,00 pada 1 jam, 90,09/94,52 pada 3 jam,
          85,40/91,35 pada 6 jam, 64,63/76,63 pada 24 jam
          [terukur, \textit{build} 2026-09-06].
    \item \textbf{NF-01 dievaluasi terhadap setiap jendela} (D-12), bukan hanya
          terhadap jendela primer. Dinyatakan \textbf{tidak terpenuhi} pada
          resolusi jam, dengan alasannya.
    \item \textbf{Jendela 6 jam adalah nominasi demi keterbacaan, bukan optimum
          yang diturunkan} --- tidak ada kriteria seleksi yang pernah
          diterapkan padanya (D-04). Nyatakan begitu.
    \item \textbf{Kalibrasi di samping setiap baris berjendela}: rasio rerata
          prediksi terhadap rerata observasi. Terukur pada tropis fold 2,
          rerata GFD prediksi 6 jam adalah \textbf{0,37 sampai 0,45} dari rerata
          observasi \textbf{di seluruh tangga} --- setiap model, sehingga ini
          sifat transformasinya dan bukan sifat modelnya (D-03)
          [terukur, tercatat].
\end{itemize}

\subsection{Hasil per Tahap Hurdle}

\begin{itemize}
    \item \textbf{Tabel: hasil tahap okurensi dan tahap cacah secara terpisah}
          (D-13), lalu cacah harapan gabungan
          $P(\text{okurensi}) \times E[\text{cacah}\mid\text{okurensi}]$.
    \item Membaca kedua tahap terpisah adalah cara melihat di mana kesalahan
          berada --- kesalahan pemeringkatan okurensi dan kesalahan intensitas
          punya obat yang berbeda.
          % TODO: isi setelah run final dibaca.
\end{itemize}

\subsection{Hasil Sweep Satu Faktor}

\begin{itemize}
    \item \textbf{Sweep \texttt{train\_rows} --- efisiensi sampel} (D-14).
          \textbf{Sudah diukur dan hasilnya tidak mendukung klaim keunggulan
          data-sedikit untuk QNN.} \texttt{train\_rows} pada 500/1.000/3.000,
          tiga seed, fold 2. QNN dikurangi NN berparameter-cocok, positif
          berarti lebih buruk [terukur, tercatat]:
          \begin{itemize}
              \item 500 baris: okurensi $+0,0237$ (sd $0,0191$); cacah
                    $+0,0054$ (sd $0,0513$);
              \item 1.000 baris: okurensi $+0,0274$ (sd $0,0142$); cacah
                    $+0,0067$ (sd $0,0287$);
              \item 3.000 baris: okurensi $+0,0269$ (sd $0,0082$); cacah
                    $+0,0208$ (sd $0,0173$).
          \end{itemize}
          Pada okurensi selisihnya melampaui simpangan baku terkumpul di setiap
          anak tangga dan \textbf{tidak menyusut} sepanjang julat data 6$\times$.
          Tesis melaporkan bahwa keunggulan \textbf{dicari dengan replikasi yang
          memadai dan tidak ditemukan} --- bukan bahwa ia tidak ada.
    \item \textbf{Sweep \texttt{hour\_encoding} --- 13/14/15 fitur} (D-02).
          % TODO: isi setelah run dibaca.
    \item \textbf{Sweep \texttt{train\_zero\_ratio} --- alami/75\%/50\%} (D-07).
          % TODO: isi setelah run dibaca.
    \item \textbf{Sweep \texttt{feature\_reduction} --- None/pca8/pca6} (D-40),
          disertai \textbf{di sisi mana pembalikan paritas} sebuah hasil berada,
          karena sirkuit dibangun pada lebar tereduksi sehingga jumlah
          parameternya ikut berubah.
          % TODO: isi setelah run dibaca.
    % TODO(open): pertanyaan terbuka 23 — EMPAT sweep OFAT dideklarasikan dan
    %   BELUM DIJALANKAN. Sebutkan secara eksplisit mana yang dijalankan dan
    %   mana yang tidak; jangan menyajikan daftar sumbu seolah semuanya berjalan.
    % TODO(open): pertanyaan terbuka 21 — sumbu `shots` DICURIGAI TIDAK BEKERJA
    %   dan BELUM DIUJI. Jangan melaporkan hasil apa pun darinya.
\end{itemize}

\subsection{Variabilitas Antar Fold}

\begin{itemize}
    \item \textbf{Tabel: hasil per fold} (D-06), fold 1--3, agar pembaca melihat
          sebaran dan bukan hanya reratanya.
    \item \textbf{Keterbatasan siklus surya}: tiga fold merentang 2022--2024,
          jauh lebih pendek daripada siklus aktivitas matahari, sehingga
          variabilitas antar fold tidak boleh dibaca sebagai variabilitas
          antar tahun secara umum.
          % TODO(open): pertanyaan terbuka 25 — apakah spektrum varians yang datar
          %   bertahan di luar satu fold satu domain? Belum diuji.
\end{itemize}

\subsection{Hasil Lintas Domain}

\begin{itemize}
    \item \textbf{Tabel: empat kombinasi latih--uji} (tropis$\rightarrow$tropis,
          tropis$\rightarrow$subtropis, subtropis$\rightarrow$subtropis,
          subtropis$\rightarrow$tropis), dengan \textit{scaler} domain
          \textbf{sumber} (D-36).
    \item \textbf{Selisih lintas domain yang tersisa dibaca sebagai temuan
          transfer}, bukan sebagai kegagalan --- dan hanya sah dibaca begitu
          karena \texttt{KX} sudah dibuang dari kedua domain (D-01) sehingga
          keduanya menyajikan 13 prediktor yang sama.
    \item \textbf{Biaya membuang \texttt{KX}} (D-01) dilaporkan terpisah:
          perbandingan terhadap studi terdahulu melemah karena prediktor yang
          dipakainya tidak tersedia di kedua domain. Nyatakan.
          % TODO: isi setelah run dibaca.
\end{itemize}

\subsection{Agregasi Spasial}

\begin{itemize}
    \item \textbf{Hasil pelengkap, bukan sumbu utama kedua} (D-05). Prediksi
          diagregasikan atas blok $2\times2$ sel pada kedua domain.
    \item \textbf{Aritmetikanya penting}: GFD adalah kerapatan, sehingga
          agregasi sel adalah $\sum(\text{cacah}) / \sum(\text{luas})$,
          \textbf{tidak pernah} rerata kerapatan per sel --- luas sel berbeda
          menurut lintang.
    \item Tropis hanya 30 sel, sehingga blok $5\times5$ nyaris seluruh domain
          dan $2\times2$ adalah batasnya.
          % TODO(open): pertanyaan terbuka 19 — dimensi lintang $\times$ bujur
          %   subtropis belum dikonfirmasi (56 sel, bukan 54 titik NASA POWER).
          %   Konfirmasi sebelum mengklaim $3\times3$ tersedia di subtropis.
\end{itemize}

\subsection{Kesetaraan Dua Jalur Gradien}

\begin{itemize}
    \item \textbf{Verifikasi \textit{forward pass}}: sirkuit PennyLane dan
          Qiskit sepakat sampai $10^{-10}$ (D-28) --- ini menjawab pertanyaan
          tentang \textbf{sirkuit}.
    \item \textbf{\textit{Run} Qiskit tereduksi} (D-39): 500 baris $\times$ 15
          epoch, satu fold, satu seed, tahap okurensi saja --- menjawab
          pertanyaan tentang \textbf{pelatihan}. Ruang lingkupnya dinyatakan
          agar tidak ada pembaca yang mengasumsikan lebih; pasangan
          \textit{hurdle} penuh adalah 21,9 jam dan tidak dijalankan.
          % TODO(open): pertanyaan terbuka 24 — D-39 berstatus DIPUTUSKAN, BELUM
          %   DIJALANKAN. Jangan menulis hasilnya sebelum ada.
\end{itemize}

\section{Pembahasan Hasil Evaluasi}

\begin{itemize}
    \item \textbf{Terhadap NF-01}: dinyatakan tidak terpenuhi pada resolusi jam
          dan dijelaskan sebagai perbedaan kategori --- NF-01 ditulis terhadap
          kerapatan bulanan, sedangkan target di sini adalah proses cacah
          94--97\% nol (D-12). Dievaluasi terhadap setiap jendela, dan pada 24
          jam pangsa nol tropis 64,63\% membuat R\textsuperscript{2} dapat
          ditafsirkan meski target tetap didominasi nol.
    \item \textbf{Terhadap NF-02}: terpenuhi secara mekanisme, \textbf{belum
          diuji ujung-ke-ujung} (D-42, pertanyaan terbuka 28). Kalimat itu
          persis, tidak lebih.
    \item \textbf{Terhadap klaim QML}: efisiensi sampel dicari dan tidak
          ditemukan (D-14); paritas parameter adalah mata uang yang lemah dan
          harus dikatakan begitu (D-34).
    \item \textbf{Terhadap studi terdahulu}: perbandingan melemah oleh
          pembuangan \texttt{KX} (D-01) dan oleh perbedaan resolusi
          (D-21) --- keduanya dinyatakan, bukan diselipkan.
    \item \textbf{Ancaman terhadap validitas, seluruhnya dinyatakan di muka:}
          komposisi CG/IC subtropis (pertanyaan terbuka 10), konvensi
          \textit{flash} lawan \textit{stroke} (pertanyaan terbuka 11), konfound
          \textit{detection efficiency} (D-27), asimetri \textit{sweep}
          (D-29, D-37), dan tujuh hiperparameter tanpa alasan tercatat
          (pertanyaan terbuka 22).
\end{itemize}
